\documentclass{article}
\usepackage{graphicx} % Required for inserting images
\usepackage{xcolor}

\usepackage{amsmath}
\usepackage{amsthm}
\usepackage{thmtools} 
\usepackage{cleveref}


% Number equations within sections (instead of chapters)
\numberwithin{equation}{section}

% Declare theorem environments anchored to sections
\declaretheorem[name=Theorem,numberwithin=section]{theorem}
\declaretheorem[name=Lemma,numberlike=theorem]{lemma}
\declaretheorem[name=Proposition,numberlike=theorem]{proposition}
\declaretheorem[name=Corollary,numberlike=theorem]{corollary}
\declaretheorem[name=Definition,numberlike=theorem]{definition}
\declaretheorem[name=Assumption,numberlike=theorem]{assumption}
\declaretheorem[name=Example,numberlike=theorem]{example}
\declaretheorem[name=Remark,numberlike=theorem]{remark}

% For restatable theorems 
\declaretheorem[name=Theorem,numberwithin=section]{restatabletheorem}

\title{Learning in two-player games between
transparent opponents}
\author{Colomban Duclaux}
\date{\today}

\usepackage[backend=biber,style=alphabetic]{biblatex} % Ensure style matches your preference
\addbibresource{../references.bib} % Link your .bib file

\newcommand{\dupoc}{\texttt{DUPOC}}
\newcommand{\cupod}{\texttt{CUPOD}}
\newcommand{\cimcic}{\texttt{CIMCIC}}
\newcommand{\dimcid}{\texttt{DIMCID}}
\begin{document}

\maketitle
These notes recap the content exposed in \cite{hutter2021learningtwoplayergamestransparent}.

\section{Abstract}
Scenario in which two RL agents repeatedly play a matrix game against each other and update their parameters after each round. Decision-making is transparent.
Two algorithms: LOLA and SOS (opponent-aware learning).

\section{Introduction}
Vanilla gradient descent learners failed to learn tit-for-tat strategies that encourage mutual cooperation in iterated prisoner's dilemma.
Important point: in one-shot PD, situations that have been studied where mutual cooperation might arise are based on one condition: mutual transparency.

Different strategies such as $\epsilon-\texttt{GroundedFairBot}$ can emerge that defect with prob $1-\epsilon$ if facing a defector. They train an agent to have such property.

\section{Learning with opponent-learning awareness}
Two algorithms based on gradient descent: \textbf{Learning with Opponent-Learning Awareness (LOLA)} and \textbf{Stable Opponent Shaping (SOS)}.
Both use gradients of both players in the gradient descent part (Taylor approximation for example).
Depending on the learning rate of each player, we have some notion of how much one is \textit{looking ahead}.

\section{Discussion}
We have seen that LOLA and SOS learners can
achieve outcomes which are drastically different from those achieved by naive
gradient descent learners.

In the prisoner’s dilemma, mutual awareness of each other’s inner workings
is helpful for achieving cooperative outcomes. This idea has existed for decades
(at least since Hofstadter (1983)), and we have now validated it in the paradigm
of gradient-based machine learning.

Chicken Game can pose some problems in the setting of mutual information, as it can favorize defection.

\section{Opinion of this paper}
Shows a computational approach of open-source gt where they try to test the theory of some known results of cooperation or defection for some agents.

Could be useful to understand how learning works in practice $\to$ Here it is basically just gradient descent on the expected payoff that is computed analytically (closed-formed).

\printbibliography
\end{document}